\begin{abstract}
Audio chord recognition (ACR) models trained on human-annotated corpora face a persistent triadic performance plateau: canonical benchmarks concentrate over $80\%$ of duration in major and minor triads, leaving rare triads ({\small\texttt{sus2}}, {\small\texttt{sus4}}, {\small\texttt{dim}}, {\small\texttt{aug}}) and upper extensions ($7\text{th}$--$13\text{th}$) acoustically sample-starved. Consequently, neural classifiers collapse into majority-class predictors, and inverse-frequency loss reweighting fails to compensate for the physical absence of diverse training exemplars. We propose a Sim-to-Real pretraining framework driven by a probabilistic composition generator that synthesizes $5{,}000$ multitrack audio tracks ($460.3\,\text{h}$) using rare-chord quota scheduling, physically constrained multi-instrument voicings, and stochastic random-mix acoustic augmentation. We pretrain ChordFormer---a 4-block Conformer encoder with six decomposed output heads---on synthetic multitracks and finetune on real audio via inverted discriminative learning rates ($\mathbf{s}_{\text{lr}} = [0.2, 0.5, 1.0]$) and floor curriculum replay ($\rho_{\min}=0.10$). Across all $11{,}436{,}357$ frames of the $1{,}217$-song benchmark, our curriculum achieves best-in-class recall on underrepresented harmonies: boosting rare triad recall from $41.67\%$ to $48.31\%$ ($+6.64\,\text{pp}$; $F_1$ from $53.11\%$ to $60.94\%$), 9th extensions from $35.33\%$ to $43.58\%$ ($+8.25\,\text{pp}$), natural 9ths by $+25.16\,\text{pp}$ ($31.64\% \!\to\! 56.80\%$), altered {\small\texttt{\#9}} by $79\times$ ($0.07\% \!\to\! 5.53\%$), and flat 13ths by $7.8\times$ ($0.84\% \!\to\! 6.58\%$). While duration-weighted MIREX WCSR shifts from $81.49\%$ to $79.59\%$ due to metric bias toward dominant triads, cross-fold variance contracts across all ten evaluation rules.
\end{abstract}
